\documentclass{ximera}

\title{Labels and refs}

\begin{document}
\begin{abstract}
\end{abstract}
\maketitle

% {{\Huge \bfseries Last Updated: \today}} \\

%\section{Basic Demo}

\section{Introduction}
\label{sec:intro}

Welcome to the demo. In \autoref{sec:math}, we will discuss some math.
We also refer to Figure~\ref{fig:demo} and Table~\ref{tab:example}.
Test scrolling to the targets with some care, e.g. is rather small windows.

\section{Mathematics Example}
\label{sec:math}

Consider the Pythagorean theorem in an \verb|equation| environment
\begin{equation}
  a^2 + b^2 = c^2   \label{eq:pythagoras}
\end{equation}
and in a normal \verb|\[ ... \]| environment
\[
  a^2 + b^2 = c^2   \tag{Pythagoras} \label{eq:tag:pythagoras}
\]
or just in a \verb|center| environment with a \verb|\hypertarget|
\begin{center}
  $a^2 + b^2 = c^2$  \hfill  \hypertarget{ht:pythagoras}{(Pythagoras (hypertarget))}
\end{center}

And consider also this numbered list:
\begin{enumerate}
  \item First item \label{itm:first}
  \item Second item \label{itm:second}
\end{enumerate}


Then we can refer to the first equation \textbf{in mathmode} using \verb|\ref| : Equation~$\ref{eq:pythagoras}$ or with \verb|\eqref|: $\eqref{eq:pythagoras}$.

We can refer to the second equation \textbf{in mathmode} using \verb|\ref| : Equation~$\text{\ref{eq:tag:pythagoras}}$ or with \verb|\eqref|: $\eqref{eq:tag:pythagoras}$.

We can refer to the third one  \verb|\hyperlink|, thus \hyperlink{ht:pythagoras}{like this}.

Note that in HTML \verb|\ref| and \verb|\eqref| without mathmode might jump to BELOW Equation~\ref{eq:pythagoras} or \eqref{eq:pythagoras}, and
\verb|\ref| might need a \verb|\text| to avoid $\ref{eq:pythagoras}$.

The first equation appears on page~\pageref{eq:pythagoras}, the second on page~\pageref{eq:tag:pythagoras}. 
Note that this does not make a lot of sens online.

We refer to item~\ref{itm:first} in the list above.

\section{Figures and Tables}

\subsection{A Demo Figure}
\label{sec:fig}

\begin{image}[0.3\textwidth]
  \centering
  \includegraphics[width=0.3\textwidth]{missionPatch.jpg} % Adjust the path to your image file
\end{image}
  \captionof{figure}{A demo image}
  \label{fig:demo}

\subsection{A Demo Table}

{
  \centering
  \begin{tabular}{|c|c|}
    \hline
    A & B \\
    \hline
    1 & 2 \\
    \hline
  \end{tabular}
  \captionof{table}{An example table}
  \label{tab:example}
}

\section{Custom Text Reference}

This document started in Section~\ref{sec:intro} and discussed math in Section~\nameref{sec:math}.


\begin{remark}

    \begin{itemize}
        \item  Links should be GLOBALLY unique, so do not use the same label in different activities.
        \item  In a PDF links behave differently depending on whether the target is in the PDF or not. 
            Test carefully and thoroughly.
    \end{itemize}
\end{remark}


\section{Intended Outcome of Test}


\section{Start of Test/Demo Area}

\begin{equation}
a^2 + b^2 = c^2    \label{eq:pyth}
\end{equation}

Facemus paginam abbastanzam longam, so that we can test the scrolling to the refs.

\lipsum[1]

\lipsum[2]

\lipsum[3]

\lipsum[4]

\hrulefill

\end{document}
